\chapter{Introducción}

\section{Contexto}
\label{contexto}
    La soledad como concepto no tiene ninguna connotación negativa, ya que cuando es buscada, puede 
ser muy placentera. El problema surge cuando esa soledad no es deseada, sino impuesta por
 aspectos externos, como pueden ser la situación social, demográfica, familiar o laboral.
  Este problema afecta principalmente a la tercera edad, debido al fallecimiento de la pareja,
   a la disminución del núcleo familiar, o al distanciamiento con los círculos sociales principales.
Según \citep{RuedaEstrada2018}, es un problema que genera una preocupación creciente cada año que 
transcurre en España, no solo a la ciudadanía, sino a profesionales, políticos, geriatras,
 psicólogos, y perfiles similares.
\par


    Esta preocupación creciente ha cristalizado institucionalmente. En 2024, la \textit{Revista Española
    de Salud Pública}, editada por el Ministerio de Sanidad, define la soledad no deseada (SND) como
    \enquote{un problema silente de nuevo orden} que afecta a un número cada vez mayor de personas en España, Europa
    y el mundo, y que conlleva consecuencias negativas tanto para la salud y el bienestar emocional como 
    para las cuentas públicas \citep{JimenezRodriguez2024}. La magnitud del fenómeno se ha podido cuantificar
    con cierta precisión gracias al trabajo del \textbf{Observatorio Estatal de la Soledad No Deseada}
    \citep{SoledadES}, impulsado por Fundación ONCE y Fundación AXA, cuyo \textit{Barómetro de la Soledad no deseada}
    en España en 2024 sitúa la prevalencia general en un 20\% de la población adulta, con un 13,5\% que padece
    soledad de carácter crónico, es decir, mantenida durante más de dos años.
\par

    Conviene matizar, sin embargo, que la soledad no deseada no se distribuye de forma homogénea por grupos de edad.
    Su prevalencia dibuja una curva en forma de «U»: es elevada en la juventud, decrece en la edad adulta y vuelve
    a ascender en los últimos tramos de la vida. Entre las personas de 65 a 74 años la prevalencia es del 14,5\%,
    pero en las mayores de 75 se eleva hasta el 20\% \citep{FundacionONCE}. Traducido a cifras absolutas, 
    se estima que cerca de tres millones de personas mayores de 65 años experimentan soledad no deseada en 
    España, un fenómeno cuyo coste para el Estado se ha estimado en torno a los 14.141 millones de euros anuales
    \citep{Senara}. Más allá del volumen, lo característico de la soledad en la vejez es su intensidad
    y persistencia: a diferencia de la juventud, suele estar vinculada a pérdidas difícilmente reversibles ---
    la viudez, las limitaciones físicas, la dificultad para tejer nuevos vínculos--- y por tanto tiende a
    cronificarse \citep{FundacionONCE}.
\par

    Esta realidad presenta, además, un marcado sesgo de género. El sentimiento de soledad es sistemáticamente 
    mayor entre las mujeres, y la diferencia se agrava con la edad: en el grupo de 65 y más años, las mujeres 
    declaran una prevalencia 7,8 puntos porcentuales superior a la de los hombres (19,85\% frente al 12\%). Este
    desequilibrio guarda relación con la mayor esperanza de vida femenina y con un patrón cultural de viudez 
    asimétrico, dado que el 82\% de las personas viudas de 70 y más años son mujeres \citep{FundacionONCE}.
\par

    A los factores demográficos y de salud descritos se suma, en el caso español, un agravante de naturaleza 
    tecnológica que afecta de manera particular a la población mayor: la brecha digital. La progresiva 
    digitalización de servicios esenciales —banca, administración, sanidad— ha convertido la competencia 
    tecnológica en un requisito de facto para la plena participación social, precisamente en el momento en que 
    se cierran los canales presenciales. El 57\% de los municipios españoles, más de 4.300, carece ya de 
    oficina bancaria, lo que afecta directamente a más de 1,5 millones de personas, en su mayoría mayores 
    y residentes en el medio rural \citep{RevistaByte}. El problema no es únicamente de infraestructura 
    —la cobertura de fibra óptica alcanzaba el 86,52\% de los hogares rurales en junio de 2024— sino, 
    sobre todo, de capacitación: más de un tercio de la población española carece de competencias digitales 
    básicas, una cifra que se dispara en los tramos de edad avanzada. La socióloga 
    Noelia Morales, de la Universidad de Salamanca, ha descrito esta confluencia de desventajas con el 
    término "multiexclusión", para subrayar que el mayor —y muy especialmente el mayor rural— afronta de 
    forma simultánea la brecha generacional, la geográfica y la digital, de modo que cada nuevo servicio 
    que se digitaliza tiende a alejarlo un escalón más de la sociedad \citep{espanol}.
\par

    En síntesis, la soledad no deseada en la tercera edad constituye en España un problema de gran escala 
    —millones de personas afectadas—, elevado coste económico —decenas de miles de millones de euros anuales— 
    e impacto sanitario contrastado, sobre el que confluye además una brecha digital que limita el acceso de 
    los propios afectados a las herramientas que podrían mitigarlo. Es en esta intersección entre necesidad 
    social y barrera tecnológica donde se justifica la propuesta del presente Trabajo Fin de Grado.
\par

\section{Motivación}

La motivación de este proyecto tiene un enfoque dual. Por un lado, en lo que al 
proyecto software se refiere, la motivación principal se basa en reforzar y 
profundizar las competencias obtenidas a lo largo de toda la formación académica 
del estudiante: desde la planificación y la toma de decisiones en un proyecto real 
hasta las tareas de ingeniería de requisitos, ingeniería del software, gestión de 
proyectos y modelado, aplicando la metodología ágil KANBAN. En este sentido, el 
proyecto persigue demostrar la capacidad de afrontar un desarrollo de principio a 
fin, justificando cada decisión técnica —desde la arquitectura hasta la elección 
del stack tecnológico— frente a las restricciones reales de un desarrollo individual 
con recursos limitados.

Por otro lado, el producto final tiene como motivación ofrecer una solución parcial, 
o al menos un paso en la dirección correcta, al problema expuesto en el contexto de 
este mismo trabajo, brindando una alternativa que estimule la mente y apacigüe esa 
sensación de soledad no deseada (SND) que sufre un porcentaje considerable de la 
sociedad española. Esta vertiente social se sustenta en una constatación que emerge 
del análisis del mercado actual: las soluciones tecnológicas existentes, pese a su 
proliferación, presentan carencias estructurales que dejan desatendido precisamente 
al perfil de usuario más vulnerable. Por un extremo, los robots sociales de gama alta 
ofrecen una experiencia cuidada pero a un coste prohibitivo para la economía del 
pensionista medio español y sin presencia ni localización al castellano. Por el 
extremo opuesto, los compañeros conversacionales de inteligencia artificial de uso 
masivo no han sido diseñados para combatir la soledad de forma terapéutica, sino para 
maximizar el tiempo de uso, lo que ha derivado en patrones documentados de dependencia 
emocional y en una creciente preocupación regulatoria a escala europea.

**NOTE: revisar asunción de que mi aplicación sí combate la soledad de forma terapéutica
, ya que no está justificado científicamente**
% Cuidado con la frase "los compañeros conversacionales de inteligencia artificial de uso 
% masivo no han sido diseñados para combatir la soledad de forma terapéutica" ya que da a 
% entender que mi aplicación sí que hace eso, y no es cierto


Surge así la motivación de explorar una tercera vía: una aplicación accesible, 
construida sobre hardware de uso común y con un coste de despliegue marginal, diseñada 
desde su concepción en torno a la accesibilidad y a la lengua española, e ideada de 
manera que la tecnología actúe como complemento del vínculo humano y no como su 
sustituto. Esta última premisa, materializada en el modelo dual que vincula a la 
persona mayor con un familiar, constituye el eje diferencial de la propuesta y la 
respuesta directa a los riesgos identificados en las alternativas comerciales. El 
diseño concreto de esta solución, así como los objetivos que persigue, se detallan 
en los siguientes apartados.


\section{Descripción de la Solución}

Frente a las carencias identificadas en las soluciones existentes (expuestas en el Estado del Arte
\ref{estado-arte}), este Trabajo Fin de Grado propone \textbf{Quiet Companion}, una aplicación
 móvil multiplataforma concebida para combatir la soledad no deseada en personas mayores mediante 
 un asistente conversacional por voz. La propuesta se articula en torno a tres principios 
de diseño que responden directamente a los problemas descritos en el contexto: la 
\textit{accesibilidad radical} de la interfaz, el uso del \textit{lenguaje natural 
hablado} como vía pre-establecida y principal de interacción, y un modelo 
de \textit{cuidado distribuido} que vincula a la persona mayor con un familiar o acompañante.

\subsection*{El núcleo conversacional}

El corazón de la aplicación es una única pantalla dominada por un botón de micrófono 
de gran tamaño y alto contraste. Al pulsarlo, el sistema transcribe la voz del usuario 
a texto (\textit{speech-to-text}), la envía a un modelo de lenguaje en la nube que 
genera una respuesta empática y adaptada, y reproduce dicha respuesta mediante síntesis 
de voz (\textit{text-to-speech}). La interacción prescinde deliberadamente de gestos 
complejos —no se emplean deslizamientos, doble toque ni pulsaciones prolongadas— y se 
limita a la pulsación simple, la modalidad de interacción más robusta frente a las 
limitaciones motoras y visuales propias de la edad. El estado del sistema 
(«Escuchando», «Pensando», «Hablando») se comunica de forma visual y con tipografía 
de gran tamaño, conforme a los criterios de perceptibilidad y operabilidad de la 
norma WCAG 2.2.

**NOTE: revisar al finalizar la implementación si he mantenido las mismas etiquetas de estado
de sistema**

\subsection*{Revelación progresiva: una respuesta al espectro de autonomía}

La aportación conceptual más relevante de la propuesta es el principio de 
\textit{revelación progresiva}, que aborda una realidad frecuentemente ignorada por 
las soluciones del mercado: la persona mayor no es ni plenamente autónoma ni plenamente 
dependiente, sino que se sitúa en un punto variable de un espectro. La aplicación no 
impone un único modo de uso, sino que se adapta a la capacidad de cada usuario.

Al abrir la aplicación por primera vez, el usuario elige entre dos modos mediante una 
pantalla con dos opciones claras y de gran tamaño:

\begin{itemize}
    \item \textbf{Modo Independiente}: dirigido a la persona mayor con capacidad y 
    voluntad de gestionar la herramienta por sí misma. En este modo se habilita un 
    acceso a los ajustes de accesibilidad en su propio dispositivo.
    \item \textbf{Modo Companion}: dirigido a la persona mayor que delega la configuración 
    en un familiar. La aplicación genera un código de vinculación corto que el familiar 
    introduce en su propio terminal; hasta que se completa el emparejamiento, la 
    interfaz del mayor permanece bloqueada en un «modo ultra-simple» que oculta los 
    ajustes para evitar manipulaciones accidentales.
\end{itemize}

Esta decisión es a la vez social y técnica. Socialmente, respeta la autonomía del 
mayor sin imponer paternalismo, dándole el control únicamente si demuestra capacidad 
para gestionarlo. Técnicamente, resulta eficiente: el núcleo de la aplicación ---la 
pantalla de conversación por voz--- es idéntico para ambos perfiles, y la diferencia 
entre modos se reduce a mostrar u ocultar condicionalmente el acceso a ajustes, lo 
que minimiza la duplicación de código y de diseño.

\subsection*{Configuración adaptativa y cuidado distribuido}

Tanto el mayor independiente como el familiar vinculado pueden ajustar los parámetros 
de accesibilidad ---volumen, velocidad de la voz y tamaño del texto--- a través de una 
interfaz simplificada. Cuando existe un familiar vinculado, los cambios que este 
realiza desde su panel remoto se propagan automáticamente al dispositivo del mayor, 
sin que este deba intervenir. Este mecanismo materializa el modelo de cuidado 
distribuido que diferencia la propuesta de los compañeros conversacionales 
generalistas: la presencia de un observador humano de confianza no solo facilita la 
configuración, sino que sienta las bases para una futura monitorización del bienestar 
emocional, garantizando que la tecnología complemente el vínculo humano en lugar de 
sustituirlo.

\section{Objetivos}

\subsection*{Objetivo general}

El objetivo general de este Trabajo Fin de Grado es \textbf{diseñar, desarrollar y 
validar un prototipo funcional (MVP) de una aplicación móvil que, mediante un asistente 
conversacional por voz accesible, contribuya a mitigar la soledad no deseada en 
personas mayores}, aplicando para ello metodologías y buenas prácticas de ingeniería 
del software propias de un entorno profesional.

\subsection*{Objetivos específicos}

Este objetivo general se descompone en los siguientes objetivos específicos, que 
constituyen el alcance comprometido del proyecto:

\begin{enumerate}
    \item \textbf{OE-1.} Implementar un sistema de gestión de perfiles con revelación 
    progresiva que permita al usuario operar en modo independiente o vinculado a un 
    familiar, así como el mecanismo de emparejamiento entre ambos. (\textit{\hyperref[tab:hu1]{HU-1}})

    \item \textbf{OE-2.} Desarrollar una interfaz conversacional por voz que integre 
    el ciclo completo de reconocimiento de voz, generación de respuesta mediante un 
    modelo de lenguaje y síntesis de voz, comunicando visualmente sus estados. 
    \textit{\hyperref[tab:hu2]{HU-2}}

    \item \textbf{OE-3.} Garantizar el cumplimiento de los criterios de accesibilidad 
    de la norma WCAG 2.2 \citep{wcag} en la interfaz dirigida a la persona mayor, mediante 
    tipografías de gran tamaño, áreas táctiles amplias e interacción simplificada. 
    \textit{\hyperref[tab:hu3]{HU-3}}

    \item \textbf{OE-4.} Implementar un sistema de configuración adaptativa que permita 
    ajustar los parámetros de accesibilidad y que propague en tiempo real los cambios 
    realizados por el familiar al dispositivo del mayor. \textit{\hyperref[tab:hu4]{HU-4}}

    \item \textbf{OE-5.} Diseñar e implementar la solución sobre una arquitectura 
    escalable y mantenible, aplicando los principios de Clean Architecture y 
    Feature-Sliced Design, así como gestionar el ciclo de vida del proyecto mediante 
    una metodología ágil basada en tableros KANBAN.
\end{enumerate}

\subsection*{Líneas de trabajo futuro}

El diseño completo del producto contempla funcionalidades adicionales que, por 
restricciones de tiempo propias de un Trabajo Fin de Grado, quedan fuera del alcance 
del presente MVP y se proponen como líneas de trabajo futuro:

\begin{itemize}
    \item Dotar al asistente de \textbf{memoria conversacional}, inyectando en el 
    contexto del modelo el nombre y los datos básicos del usuario. \textit{(HU-5)}

    \item Incorporar un módulo de \textbf{seguimiento del estado de ánimo} que catalogue 
    el sentimiento de las conversaciones y lo presente al familiar como indicador 
    agregado de bienestar. \textit{(HU-6)}

    \item Añadir una función de \textbf{llamada de emergencia} accesible desde la 
    pantalla principal.
\end{itemize}

\section{Estructura de la memoria}

Esta memoria se ha dividido en varios capítulos estructurados. Para un seguimiento y comprensión total 
del trabajo realizado, se proporciona una breve explicación del contenido:

\begin{itemize}
    \item El \textbf{capítulo 1} detalla el ámbito del trabajo, explicando el contexto, motivación 
    y objetivos de la creación de la aplicación.
    \item El \textbf{capítulo 2} ofrece una visión global del estado del arte y la oferta actual de 
    soluciones similares en el mercado.
    \item El \textbf{capítulo 3} analiza las metodologías de desarrollo consideradas para este proyecto,
     exponiendo la aplicación de la metodología elegida, KANBAN, justificando su idoneidad para este proyecto. 
    \item El \textbf{capítulo 4} detalla los actores del sistema, los requisitos funcionales y no funcionales, 
    y los casos de uso, culminando con el modelo conceptual de la aplicación.
    \item El \textbf{capítulo 5} muestra la fase de diseño y desarrollo del sistema. % Habría que separarlo en dos?
    \item El \textbf{capítulo 6} explica las diferentes validaciones y pruebas que se han efectuado sobre el software.
    \item El \textbf{capítulo 7} muestra la planificación y realiza un estudio de costes del proyecto en su totalidad
    \item El \textbf{último capítulo} expone las conclusiones del proyecto, detallando los objetivos cumplidos 
    , las limitaciones finales, y las líneas de trabajo futuras.
\end{itemize}

Finalmente, se incluyen las referencias bibliográficas utilizadas.

La aplicación resultante, el código fuente y el manual se pueden encontrar en el repositorio:  
\url{https://github.com/maario106/quiet_companion}